\scanpage{102}
\mspage{\unsure{89|109}}
% archive sheet number: 101

i.e. une catégorie avec $W,C$ \deleted{\unsure{fibr. stable à}, $C$ stable par} \inserted{satisfaisant}
$\mathrm{Cof}_{0},\mathrm{Cof}_{1},\mathrm{Cof}_{2}$ de la page 39, \unsure{plus} $TF\overset{\mathrm{déf}}=C_{*}\subset W$,
plus l'axiome de factorisation pour $(C,TF)$
(lequel implique $\mathrm{Cof}_{3}$).

\uline{Corollaire} 5.5. Soit \deleted{\unsure{$C$,}} \inserted{$(M,W,C)$} une catégorie à cofibrations
spéciale, telle qu'il existe une \uline{petite}
partie $TC_{0}\subset TC=C\cap W$, avec
\[
TC\subset {TC_{0}}^{\sim}, \tag*{(5.5.1)}
\]
ou ce qui revient au $=$, $(TC)^{\sim}=(TC_{0})^{\sim}$ i.e.
\[
\underset{\substack{\Vert\,\mathrm{déf}\\(TC)_{*}}}{F}=(TC_{0})_{*}.\tag*{(5.5.2)}
\]
Alors \deleted{\unsure{$(W,C$|$(M,W,C$}} $(W,C,F)$ est un triple de
Quillen. \deleted{clos ssi} Ce triple est clos ssi $C$ est
stable par facteurs directs.

En effet, le th. de factorisation de
\uline{XV} \S~1 \unsure{implique|infère} que moyennant l'hypo \unsure{faite|faible}
\deleted{\unsure{sur}} $TC_{0}\subset TC$, le couple $(TC,F)$ \deleted{\unsure{satisfait}}
satisfait à la condition de factorisation.
En effet, par les conditions $b,b'$ de la
définition 4.4., $TC$ a les propriétés de
stabilité voulues (\unsure{contient} $\unsure{\mathrm{iz}|\mathrm{is}}$, stable par
\unsure{cochgt|cochgt.} de base et par composition transfinie).

